\chapter*{Chapitre 4 — Production de chaleur}
\addcontentsline{toc}{chapter}{Chapitre 4 — Production de chaleur}

% ============================================================
\section*{Énoncés}
% ============================================================

\exercice{Dimensionnement de la puissance d'une chaudière collective}

Un immeuble collectif de 20 logements de 70\,m² chacun est situé en zone climatique H2 (température de base $T_{base} = -7$\,°C, $T_{int} = 19$\,°C). Le coefficient de déperdition moyen est $G = \SI{1{,}2}{\watt\per\square\metre\per\kelvin}$. Les besoins ECS représentent \SI{8}{\kilo\watt} en pointe. Le coefficient de foisonnement est 0,8.

\begin{enumerate}
  \item Calculer les déperditions totales de l'immeuble.
  \item Calculer la puissance de production nécessaire (chauffage + ECS).
  \item Choisir entre deux chaudières à condensation : modèle A (\SI{60}{\kilo\watt}) et modèle B (\SI{80}{\kilo\watt}). Justifier.
\end{enumerate}

% ----------------------------------------------------------

\exercice{Comparaison PAC vs chaudière à condensation en rénovation}

Un propriétaire doit remplacer sa vieille chaudière gaz dans une maison de 120\,m² (zone H1, $T_{base} = -12$\,°C). Les déperditions sont estimées à \SI{10}{\kilo\watt}. Les émetteurs sont des radiateurs haute température (80/60\,°C). Il hésite entre :
\begin{itemize}
  \item Solution A : chaudière à condensation gaz, $\eta_{PCS} = 104\,\%$ (référencé PCS), prix gaz \SI{0{,}11}{\euro\per\kilo\watt\heure}
  \item Solution B : PAC air/eau haute température, SCOP = 2,8 (fonctionnement à 60\,°C), prix élec \SI{0{,}22}{\euro\per\kilo\watt\heure}
\end{itemize}
Consommation annuelle estimée : 15\,000 kWh thermiques.

\begin{enumerate}
  \item Calculer le coût annuel de chauffage pour chaque solution.
  \item Sachant que la PAC coûte \SI{5000}{\euro} de plus à l'installation, calculer le temps de retour sur investissement simple.
\end{enumerate}

% ----------------------------------------------------------

\exercice{Analyse des performances d'une PAC en exploitation}

Une PAC air/eau est relevée lors de sa mise en service par les mesures suivantes :
\begin{itemize}
  \item Puissance électrique mesurée : \SI{3{,}8}{\kilo\watt}
  \item Température départ eau : \SI{45}{\celsius}
  \item Température retour eau : \SI{40}{\celsius}
  \item Débit eau : \SI{1200}{\litre\per\hour}
  \item Température extérieure : \SI{2}{\celsius}
\end{itemize}
$\rho_{eau} = \SI{992}{\kilogram\per\cubic\metre}$ à 42\,°C, $c_p = \SI{4179}{\joule\per\kilogram\per\kelvin}$.

\begin{enumerate}
  \item Calculer la puissance thermique produite.
  \item Calculer le COP mesuré.
  \item Le fabricant annonce COP = 3,8 en A2/W45. Analyser l'écart.
\end{enumerate}

\newpage
% ============================================================
\section*{Corrigés}
% ============================================================

\subsection*{Corrigé — Exercice 1}

\textit{Question 1 — Déperditions :}
\[
  S_{totale} = 20 \times 70 = \SI{1400}{\square\metre}
\]
\[
  \dot{Q}_{dep} = G \times S \times (T_{int} - T_{base}) = 1{,}2 \times 1400 \times (19 - (-7)) = 1{,}2 \times 1400 \times 26 = \SI{43\,680}{\watt} \approx \SI{43{,}7}{\kilo\watt}
\]

\[
  \boxed{\dot{Q}_{dep} \approx \SI{43{,}7}{\kilo\watt}}
\]

\textit{Question 2 — Puissance nécessaire :}
\[
  \dot{Q}_{total} = (\dot{Q}_{dep} + \dot{Q}_{ECS}) \times \frac{1}{foisonnement} = (43{,}7 + 8) \times \frac{1}{0{,}8} = 51{,}7 \times 1{,}25 = \SI{64{,}6}{\kilo\watt}
\]

\[
  \boxed{\dot{Q}_{total} = \SI{64{,}6}{\kilo\watt}}
\]

\textit{Question 3 — Choix :} Le modèle A (\SI{60}{\kilo\watt}) est insuffisant (60 < 64,6). Le \textbf{modèle B (\SI{80}{\kilo\watt})} est retenu — il couvre les besoins avec une marge de 24\,\%, suffisante sans surdimensionnement excessif.

% ----------------------------------------------------------

\subsection*{Corrigé — Exercice 2}

\textit{Question 1 — Coûts annuels :}

\textit{Solution A (chaudière) :} L'énergie achetée en gaz pour produire 15\,000\,kWh thermiques :
\[
  E_{gaz} = \frac{15\,000}{1{,}04} = \SI{14\,423}{\kilo\watt\heure_{PCS}}
\]
\[
  C_A = 14\,423 \times 0{,}11 = \SI{1587}{\euro\per an}
\]

\textit{Solution B (PAC) :} L'énergie électrique consommée :
\[
  E_{élec} = \frac{15\,000}{2{,}8} = \SI{5357}{\kilo\watt\heure_{élec}}
\]
\[
  C_B = 5357 \times 0{,}22 = \SI{1179}{\euro\per an}
\]

\[
  \boxed{C_A = \SI{1587}{\euro\per an} \quad \text{vs} \quad C_B = \SI{1179}{\euro\per an}}
\]

\textit{Question 2 — Temps de retour :}
\[
  \text{Économie annuelle} = 1587 - 1179 = \SI{408}{\euro\per an}
\]
\[
  TRI = \frac{5000}{408} \approx \mathbf{12{,}3\,\text{ans}}
\]

\[
  \boxed{TRI \approx 12{,}3 \text{ ans}}
\]

\begin{attention}
Ce temps de retour est à comparer avec les aides disponibles (MaPrimeRénov', CEE) qui peuvent le réduire significativement.
\end{attention}

% ----------------------------------------------------------

\subsection*{Corrigé — Exercice 3}

\textit{Question 1 — Puissance thermique :}
\[
  \dot{m} = \frac{1200 \times 992}{3600 \times 1000} = \SI{0{,}331}{\kilogram\per\second}
\]
\[
  \dot{Q} = \dot{m} \cdot c_p \cdot \Delta T = 0{,}331 \times 4179 \times 5 = \SI{6\,916}{\watt} \approx \SI{6{,}9}{\kilo\watt}
\]

\[
  \boxed{\dot{Q} \approx \SI{6{,}9}{\kilo\watt}}
\]

\textit{Question 2 — COP mesuré :}
\[
  COP_{mes} = \frac{6{,}9}{3{,}8} = \mathbf{1{,}82}
\]

\[
  \boxed{COP_{mes} = 1{,}82}
\]

\textit{Question 3 — Analyse :} Le COP mesuré (1,82) est très inférieur au COP annoncé (3,8 en A2/W45). Pistes de diagnostic :
\begin{itemize}
  \item Vérifier l'état du circuit frigorifique (charge en fluide frigorigène)
  \item Vérifier l'encrassement ou le givrage de l'échangeur air/fluide
  \item Vérifier que la PAC n'est pas en mode dégivrage lors de la mesure
  \item Vérifier le débit d'air sur l'évaporateur (filtre encrassé ?)
  \item Vérifier l'absence de court-circuit d'air chaud recyclé vers l'aspiration
\end{itemize}
